\subsection{Covariance-aware integer portfolios}
\label{sec:ellipsoid}

The threshold rules in Section~\ref{sec:swaps} restrict which pairs of countries
can be exchanged. We now allow any portfolio at the same quota and model the joint
uncertainty in its estimated gain. This matters when the same fitted heat-sensitivity
model drives the estimates for several countries. Adding two countries with similarly
uncertain slopes may add risk, while exchanging two countries with similar
sensitivities to estimation error can cancel that error.

\subsubsection{Pooling slopes and estimating joint uncertainty}

Let $G$ be the matrix of country-level training scores. Its three columns use
the score specifications from Section~\ref{sec:swaps}. For training country $i$,
the feature basis $\phi_i$ contains an intercept, standardised covariates,
the transformation $\tanh(2z)$ of standardised covariates, and pairwise products
of the transformed covariates. We standardise using the training set. With
$\Phi=(\phi_i^\top)_i$, we estimate all three surfaces by
\[
 \widehat\Theta=(\Phi^\top\Phi+10P)^{-1}\Phi^\top G,
 \qquad P=\mathrm{diag}(0,1,\ldots,1).
\]
The ridge penalty pools estimates towards a common slope while allowing
smooth nonlinearities and interactions. Let $\widehat\theta$ average the
three coefficient vectors and let $H$ contain the feature basis for the
assignment batch. Define $C=(\Phi^\top\Phi+10P)^{-1}\Phi^\top$, the residual
$\widehat e=G\mathbf{1}/3-\Phi\widehat\theta$, and the ridge leverage
$h_i=(\Phi C)_{ii}$. For country $i$, a leverage-adjusted residual
$\widehat e_i/\max(1-h_i,0.05)$ scales its influence on the fitted slopes.
We use the matrix
\[
 A=\left[
   HC\,\mathrm{diag}\!\left(\frac{\widehat e_i}{\max(1-h_i,0.05)}\right)
   \;\middle|\;
   \frac{H(\widehat\Theta-\widehat\theta\mathbf{1}^\top)}{\sqrt{3}}
   \right].
\]
The first block is a ridge sandwich factor and the second encodes
disagreement across the fitted surfaces. These blocks can imply correlated
errors for countries with similar covariates. This construction is a working
uncertainty model. It does not establish confidence coverage or capture
unmeasured confounding.

Let $D$ be the diagonal matrix of population-scaled, equity-weighted
exposure in the assignment batch. Country benefit has centre
$\mu=DH\widehat\theta$ and uncertainty factor $L=DA$. The ambiguity set
contains benefit vectors $\mu+Lu$ with $\|u\|_2\leq r$. We fix the radius
at $r=1.645$ in the benchmark. It acts as a risk penalty and has no
guaranteed probability interpretation.

\subsubsection{Conic formulation and certificate}

A portfolio is a binary vector $x$ with $\mathbf{1}^\top x=q$, and the reference
$x_0$ has the same number of places. Write $d=x-x_0$. Cauchy--Schwarz gives the
worst paired gain in the ellipsoid, including the case $L^\top d=0$:
\begin{equation}
 \min_{\|u\|_2\leq r} d^\top(\mu+Lu)
 = \mu^\top d-r\|L^\top d\|_2 .
 \label{eq:ellipsoid}
\end{equation}
For $L^\top d\ne0$ the minimum is attained by
$u=-rL^\top d/\|L^\top d\|_2$. Maximising \eqref{eq:ellipsoid} over the quota
constraint is a mixed-integer second-order cone programme \citep{bentalnemirovski1998}. Both the signs and
off-diagonal entries of $LL^\top$ matter, because $d$ is positive for additions
and negative for removals. With radius zero the solution selects the $q$
largest fitted benefits. With positive radius the ranking is a portfolio
decision because each country's contribution to risk depends on the other
additions and removals.

The implementation uses a linear integer master problem with an objective
variable $t$ and a progressively enlarged set of cuts. For every $\|v\|_2\leq1$,
the inequality $t\leq(\mu-rLv)^\top d$ is valid. At a master solution that
overstates the worst-case gain, choosing
$v=L^\top d/\|L^\top d\|_2$ gives the tightest cut at that portfolio. The
master dual bound is an upper bound on the conic optimum. The best feasible
portfolio evaluated directly in \eqref{eq:ellipsoid} supplies a lower bound.
We initialise it with the reference and pairwise improvements from the
reference and the zero-radius optimum. The difference between the bounds
certifies the remaining optimisation error, up to numerical solver tolerance.
The certificate applies to the fitted objective. Its scope does not extend to
the realised policy gain.

\subsubsection{Benchmark design}

We compare full-quota optimisation with zero-radius pooling, robust threshold
swaps and the exposure reference. The new simulations draw 150 training
countries and an independent batch of 60 test countries in the three designs
from Section~\ref{sec:swaps}. Each method chooses 18 of the 60 countries,
and the oracle ranks the 60 countries by their true benefits. In the agricultural
application we use the same 126-country sample, nested outcome scoring and
exactly 37 places as in Section~\ref{sec:swaps}. Nuisance fitting and pooling
exclude all outcome data from the batch being assigned. We evaluate every
portfolio on the same held-out DR scores and rerun evaluation with alternative
clipping. Intervals are conditional on the estimated scores and fitted rules.
